\subsubsection{Gaussian Naive Bayes vs. LDA}
\paragraph{Lý thuyết}
Gaussian Naive Bayes (GNB) và Linear Discriminant Analysis (LDA) đều là các mô
hình phân loại dựa trên xác suất, nhưng khác nhau ở mức độ giả thiết về phân
phối dữ liệu.

Gaussian Naive Bayes giả định rằng các đặc trưng độc lập có điều kiện theo
lớp:
\begin{equation}
  p(\boldsymbol{x} \mid t=k)
  = \prod_{j=1}^{d} p(x_j \mid t=k),
\end{equation}
với mỗi đặc trưng tuân theo phân phối Gaussian:
\begin{equation}
  p(x_j \mid t=k)
  = \mathcal{N}(\mu_{kj}, \sigma_{kj}^2).
\end{equation}

Do giả thiết độc lập, ma trận hiệp phương sai của mỗi lớp là ma trận đường
chéo, không mô hình hóa tương quan giữa các đặc trưng.

Ngược lại, LDA giả định rằng dữ liệu mỗi lớp tuân theo phân phối Gaussian đa
biến với cùng ma trận hiệp phương sai:
\begin{equation}
  p(\boldsymbol{x} \mid t=k)
  = \mathcal{N}(\boldsymbol{\mu}_k, \boldsymbol{\Sigma}).
\end{equation}

Điều này cho phép LDA mô hình hóa mối tương quan giữa các đặc trưng, vì
$\boldsymbol{\Sigma}$ không bị ràng buộc là đường chéo.

Hàm phân biệt của LDA có dạng:
\begin{equation}
  \delta_k(\boldsymbol{x}) =
  \boldsymbol{x}^\intercal \boldsymbol{\Sigma}^{-1} \boldsymbol{\mu}_k
  - \frac{1}{2}
  \boldsymbol{\mu}_k^\intercal \boldsymbol{\Sigma}^{-1}
  \boldsymbol{\mu}_k
  + \log \pi_k.
\end{equation}

Từ đó, biên quyết định của LDA là tuyến tính, trong khi GNB tạo biên phi tuyến
do giả thiết độc lập từng chiều.

\paragraph{Thực nghiệm}
Trong thực tế, hiệu năng của Gaussian Naive Bayes và LDA phụ thuộc mạnh vào
đặc điểm dữ liệu.

Gaussian Naive Bayes thường hoạt động tốt hơn LDA trong các trường hợp:
\begin{itemize}
  \item Số lượng mẫu nhỏ so với số chiều dữ liệu
  \item Các đặc trưng gần như độc lập hoặc đã được tiền xử lý
  \item Dữ liệu dạng sparse như văn bản (bag-of-words, TF-IDF)
\end{itemize}

Trong các trường hợp này, việc ước lượng ma trận hiệp phương sai trong LDA trở
nên không ổn định, trong khi GNB chỉ cần ước lượng trung bình và phương sai
từng chiều.

Ngược lại, LDA thường cho kết quả tốt hơn khi:
\begin{itemize}
  \item Các đặc trưng có tương quan đáng kể
  \item Số lượng mẫu đủ lớn để ước lượng tốt
  \item Cấu trúc dữ liệu gần với giả định Gaussian đa biến
\end{itemize}

Tóm lại, GNB có giả thiết mạnh hơn nhưng ít tham số hơn, giúp ổn định trong
các bài toán dữ liệu nhỏ hoặc nhiều chiều. Trong khi đó, LDA tổng quát hơn và
khai thác tốt mối quan hệ giữa các đặc trưng khi dữ liệu đủ lớn và phù hợp giả
định.
